% The "Beyond TD3: other directions" slide that used to follow has been removed.
% It listed NAF and Input Convex Neural Networks -- two dead ends students will
% never meet again -- next to two cropped screenshot quotes from the SAC paper.
% SAC and HER are both already pointed at from What's Next, where they belong.
\section{TD3: Twin Delayed DDPG}
\begin{frame}{TD3: DDPG Done Right \rref{4}}
\vspace{0.3em}
The experiments showed DDPG \emph{overestimates} $Q$ and is brittle. \textbf{TD3} (Twin Delayed DDPG) keeps the deterministic off-policy actor-critic and adds three fixes:
\begin{enumerate}
    \footnotesize
    \setlength{\itemsep}{0.5em}
    \item \textbf{Clipped double-Q (twin critics):} learn two critics $Q_{\phi_1}, Q_{\phi_2}$ and use the \emph{smaller} of the two in the TD target, so the actor cannot exploit whichever one is optimistic:
    \[
    y = r + \gamma \min_{i=1,2} Q_{\phi_i^{-}}\!\big(s',\, \tilde a\big)
    \]
    \item \textbf{Target policy smoothing:} $\tilde a = \mathrm{clip}\big(\mu_{\theta^{-}}(s') + \epsilon,\; a_{\min},\, a_{\max}\big)$ with $\epsilon \sim \mathcal{N}(0,\sigma)$ clipped to $[-c,c]$ --- blurs the target so a sharp, spurious peak in $Q$ cannot be chased.
    \item \textbf{Delayed policy updates:} update the actor and the targets less often than the critics (e.g.\ once every 2 critic steps), so the actor climbs a critic that has had time to settle.
\end{enumerate}
\vspace{0.4em}
\footnotesize\textcolor{red}{Takeaway:} TD3 is the deterministic successor to DDPG, and its clipped double-Q idea carries directly into SAC (Day 6).
\end{frame}
